\subsection{Inconsistency Measurement}

Inconsistency measurement provides formal methods for quantifying the degree of inconsistency within logical knowledge bases. Rather than treating inconsistency as a binary property, these measures assign numerical values that reflect the severity or extent of conflicts, enabling comparative analysis and systematic approaches to conflict resolution~\cite{Thimm2019}.

\subsubsection{Inconsistency Measures and Rationality Postulates}

An inconsistency measure provides a systematic way to assign numerical values reflecting the degree of inconsistency.

\begin{defn}[Inconsistency Measure]
    \label{defn:inconsistency_measure}
    An \emph{inconsistency measure} is a function $I : \mathcal{K} \to \mathbb{R}_{\geq 0}$ that maps every knowledge base from a given class $\mathcal{K}$ to a non-negative real number.
\end{defn}

Inconsistency measures are typically required to satisfy certain rationality postulates that capture intuitive properties of inconsistency quantification~\cite{Hunter2008,Thimm2019}:

\begin{enumerate}
    \item \textbf{Consistency:} $I(K) = 0$ if and only if $K$ is consistent.
    \item \textbf{Monotonicity:} If $K \subseteq K'$, then $I(K) \leq I(K')$.
    \item \textbf{Free Formula Independence:} $I(K \cup \{\phi\}) = I(K)$ if $\phi$ is not involved in any inconsistency of $K \cup \{\phi\}$.
\end{enumerate}

These postulates may require modification when adapting measures to specific domains. In particular, Monotonicity assumes that adding information can only increase inconsistency, which does not hold in non-monotonic contexts.

\subsubsection{The Contension Measure}

Several approaches to inconsistency quantification exist, including syntactic approaches based on minimal inconsistent subsets and semantic approaches based on model-theoretic analysis~\cite{Thimm2019}. Among semantic approaches, the contension measure $I_c$~\cite{Grant2011} is particularly relevant as a baseline for domain adaptations.

The contension measure employs three-valued paraconsistent semantics~\cite{Priest1979}, where propositions can be assigned a third truth value $B$ (``both true and false'') representing conflict. The measure counts the minimum number of propositions that must be assigned $B$ to satisfy the knowledge base:
\begin{equation*}
    I_c(K) = \min\{|\upsilon^{-1}(B) \cap \text{At}| \mid \upsilon \models_3 K\}
\end{equation*}

This proposition-level analysis provides finer granularity than formula-based measures: while $\{a, \neg a\}$ and $\{a \land b \land c, \neg a \land \neg b \land \neg c\}$ both contain one minimal inconsistent subset, $I_c$ yields 1 and 3 respectively, reflecting the number of conflicting propositions~\cite{Thimm2019}. Section~\ref{sec:naive_encoding} demonstrates that despite this granularity, $I_c$ yields uninformative results when applied to naive propositional encodings of planning problems.

\subsubsection{Domain Adaptation Precedents}

Inconsistency measurement techniques have been successfully adapted to diverse logical domains, establishing methodological precedent for this thesis.

\cite{Ulbricht2016} adapted inconsistency measures to \ac{ASP}, addressing the fundamental challenge that non-monotonic reasoning violates classical monotonicity postulates. Their key observation, that adding information to a non-monotonic knowledge base can reduce inconsistency, directly parallels the planning context, where adding operators can enable new solution paths. This motivates the Operator Monotonicity postulate developed in Section~\ref{sec:postulates}.

\cite{Corea2022} developed inconsistency measures for temporal logic applied to declarative process specifications. Their work demonstrates that traditional measures fail to adequately capture temporal operators, requiring domain-specific semantics for meaningful conflict quantification. This parallels the finding in this thesis that naive propositional encodings of planning problems yield uninformative measure values, motivating planning-native measures.

These adaptations establish that inconsistency measures can be meaningfully extended to diverse formalisms when domain-specific conflict semantics are carefully considered.
